% ===============================================================
\chapter{Kripke, `A Puzzle about Belief'}
% ===============================================================

Kripke decided to defend direct reference theory, considering he hated descriptivism. So he chose to defend it from arguments like the Strong Disquotation Argument.

\section{The Strong Disquotation Argument against direct reference theory}

\begin{definition}[Strong disquotation]
A normal speaker of language $L$ sincerely assents to a sentence $S$ in $L$ if and only if they believe what's expressed by $S$ in $L$.
\end{definition}

\begin{argument}[The Strong Disquotation Argument against direct reference]\leavevmode
\begin{premises}
  \item[(P1)] It is possible to (rationally) assent to `$\alpha$ is $\alpha$' but not `$\alpha$ is $\beta$', where $\alpha = \beta$.
  \item[(C1)] So, one can rationally believe that $\alpha$ is $\alpha$ but not that $\alpha$ is $\beta$, though $\alpha = \beta$.
  \item[(P2)] If $A$ believes that $\alpha$ is $\alpha$ but not that $\alpha$ is $\beta$, then `$A$ believes that $\alpha$ is $\alpha$' is true while `$A$ believes that $\alpha$ is $\beta$' is false.
  \item[(C2)] So, it is possible for `$A$ believes that $\alpha$ is $\alpha$' to be true while `$A$ believes that $\alpha$ is $\beta$' is false, even though $A$ is rational and $\alpha = \beta$.
  \item[(P3)] If the meaning of a name were just its referent, then any two sentences differing only in the occurrence of a co-referential name would be synonymous.
  \item[(P4)] Synonymous sentences cannot diverge in truth-values.
  \item[(C3)] Hence, the meaning of a name cannot be just its referent.
\end{premises}
\end{argument}

\section{Kripke's defence against the Strong Disquotation Argument}

If we can reproduce the effect of the Substitution Argument using disquotation instead of substitutivity, then this suggests that the root of the problem is with disquotation. So, if disquotation is independently problematic, then the problem is not with DRT.

There are two further principles to rely on to reject Strong Disquotation:

\begin{definition}\leavevmode
\begin{enumerate}[label=\roman*.]
  \item \emph{Weak disquotation}: if a normal speaker of language $L$ sincerely assents to a sentence $S$ in $L$, then they believe what's expressed by $S$ in $L$.
  \item \emph{Translation}: if a sentence $S$ in language $L$ expresses a truth in $L$, then any content-preserving translation of $S$ into $L'$ expresses a truth in $L'$.
\end{enumerate}
\end{definition}

\begin{note}
SD $\Rightarrow$ WD; so rejecting WD is enough to reject Strong Disquotation.
\end{note}

The Pierre case combines both.

\section{The Pierre case}

\begin{example}[Setup]
Pierre grows up in France, speaking only French. Based on postcards and what he has heard, he assents to `Londres est jolie'. Later, Pierre moves to an ugly part of London and learns to speak English by immersion, without ever leaving this part of the city. Based on his experience, he eventually assents to `London is not pretty'. Despite this, he never recognizes that `Londres' and `London' are names for the same city. If asked, he would assent to both sentences while speaking English in London.
\end{example}

\begin{argument}[Pierre puzzle (the actual Frege puzzle)]\leavevmode
\begin{premises}
  \item[(P1)] Pierre sincerely assents to `Londres est jolie' and `London is not pretty'.
  \item[(P2)] `Londres est jolie' translates as `London is pretty' in English.
  \item[(C1)] So, Pierre believes that London is pretty and that London is not pretty.
  \item[(P3)] That London is pretty and London is not pretty is a contradiction.
  \item[(C2)] Hence, Pierre has contradictory beliefs.
  \item[(P4)] Having contradictory beliefs is irrational.
  \item[(P5)] Being irrational requires some kind of mistake in reasoning.
  \item[(C3)] Thus, Pierre has made a mistake in reasoning.
  \item[(P6)] Pierre has made no mistake in reasoning, and his beliefs are perfectly rational.
\end{premises}
\end{argument}

\begin{enumerate}
  \item This is a puzzle because at no point did Pierre actually make a mistake in his reasoning, relative to Pierre's perspective, which contradicts C3), C2) and C1).
  \item From here, we reject both weak and strong disquotation, because despite sincerely assenting, Pierre did not believe what was expressed by the sentence he said.
\end{enumerate}

\section{The Paderewski case}

The earlier argument might rely on the translation principle. Let's try something that relies only on weak disquotation.

\begin{example}[Paderewski's double life]
Peter is familiar with the famous Polish pianist Paderewski. Based on this, Peter assents to `Paderewski has musical talent'. Later, Peter comes to learn about a Polish politician named `Paderewski', and he is generally skeptical of the musical abilities of politicians. Based on this, he assents to `Paderewski does not have musical talent'. However, Peter fails to realize that these are the same person.
\end{example}

\begin{enumerate}[label=\roman*.]
  \item By two applications of Weak Disquotation, Peter believes that Paderewski has musical talent and that Paderewski does not have musical talent.
  \item Although he may have a false belief about politicians, Peter has made no logical mistake. That is, while `Paderewski is Paderewski' is a logical truth, Peter cannot infer a contradiction without realizing that both occurrences of `Paderewski' refer to the same person, which is independent of logic.
\end{enumerate}

\section{The final argument against disquotation}

\begin{argument}\leavevmode
\begin{premises}
  \item[(P1)] If Weak Disquotation were correct, then we would wrongly charge speakers with having irrational beliefs, despite their making no logical mistakes.
  \item[(P2)] All else being equal, we ought to reject principles which require wrongly charging speakers with having irrational beliefs, despite their making no logical mistakes.
  \item[(C)] So, we ought to reject Weak Disquotation.
\end{premises}
\end{argument}

\section{Objections to the implications of Kripke's puzzles}

\begin{objection}[Limited scope]
Even if the right conclusion to draw from Kripke's cases is that substitution failures involving `believes' or `thinks' are not problematic for DRT, these are not the only cases of substitution failure Frege originally considered. There are also cases involving `knows', for instance, and also failure of substitution in direct quotation. Unless something similar can be said for these sorts of cases, this will not constitute a complete defence of Direct Reference Theory.
\end{objection}

\begin{objection}[Rational contradictory beliefs]
There is a different way of responding to Kripke's cases which does not require rejecting disquotation. Suppose we accept disquotation and consider again (P4):
\begin{premises}
  \item[(P4)] Having contradictory beliefs is irrational.
\end{premises}
This may be false. Perhaps not all contradictory beliefs are irrational. In the cases of Pierre and Peter, their contradictory beliefs are not irrational because they do not assent to the identity statements `Londres is London' and `Paderewski$_1$ is Paderewski$_2$'. So we can accept that their beliefs are contradictory without charging them with irrationality.
\end{objection}
